\chapter{SQL Level 2: Aggregation, GROUP BY and Joins}
\companion{Alex The Analyst: SQL Joins, GROUP BY and HAVING (beginner series)}{\yts{alex+the+analyst+sql+joins+group+by+having+beginner}}

Level 1 read one table row by row. Real questions are about \emph{groups} (``offer rate per job'') and about \emph{several tables} (``which company made
each offer''). This chapter adds the two tools for that: aggregation and joins. Both have exact rules for how many rows come out; most interview traps
are about getting that row count wrong.

\section{Aggregate functions: many rows in, one number out}
\stuck{SQL aggregate functions COUNT SUM AVG MIN MAX}{\yts{sql+aggregate+functions+count+sum+avg+min+max+tutorial}}
An aggregate function collapses a column of many values into one: \code{COUNT}, \code{SUM}, \code{AVG}, \code{MIN}, \code{MAX} (most databases also
have \code{STDDEV}, \code{VARIANCE}; PostgreSQL has \code{PERCENTILE\_CONT}). The rule to remember: \textbf{every aggregate ignores NULL, except
\code{COUNT(*)}}.
\runsql{L2_agg_basic}
Read it carefully. There are 11 employees but only 10 known salaries. \code{AVG} = $530/10 = 53$ (Zoya ignored). If you wanted unknown salaries to count
as zero (rarely sensible), you would get $530/11 = 48.18$. Two analysts who ``computed the average salary'' can disagree for this reason alone.
\runsql{L2_count_distinct}
23 application rows, but only 8 different candidates and 9 different jobs. ``How many candidates applied?'' is 8, not 23: the difference between
\code{COUNT(*)} and \code{COUNT(DISTINCT ...)} is the difference between events and entities.

\section{GROUP BY: one result row per group}
\stuck{SQL GROUP BY explained with examples}{\yts{sql+group+by+explained+with+examples}}
\code{GROUP BY title} splits the rows into piles, one pile per distinct title, then computes the aggregates \emph{inside each pile}. It is pandas'
\code{df.groupby('title').agg(...)}: split, apply, combine.
\runsql{L2_groupby}
\begin{example}[: do it by hand]
Jobs by title: ML Engineer max CTCs \{35, 45, 22\} $\Rightarrow$ count 3, average $102/3 = 34$, max 45. Data Scientist \{20, 25, 40, 18\} $\Rightarrow$
4, $103/4 = 25.75$, 40. Data Analyst \{12, 8, 10\} $\Rightarrow$ 3, 10, 12. That is exactly the output.
\end{example}
Grouping by two columns makes one group per \emph{combination}:
\runsql{L2_groupby_two}

\begin{trap}[: the golden rule of GROUP BY]
Every column in SELECT must either be in GROUP BY or be inside an aggregate. \code{SELECT dept, name, MAX(salary) FROM employees GROUP BY dept} is
ambiguous: a department has several names, which one should be shown? PostgreSQL and SQL Server raise an error. SQLite and old MySQL silently return
\emph{some} name (SQLite happens to return the name on the max row, MySQL an arbitrary one). Never rely on it; to get ``the employee with the top salary per
department'' use a window function (Chapter 47).
\end{trap}

\section{HAVING: filtering groups}
\stuck{SQL HAVING vs WHERE difference}{\yts{sql+having+vs+where+difference+explained}}
WHERE filters \emph{rows before} grouping; HAVING filters \emph{groups after} aggregation (recall the logical order: WHERE is step 2, HAVING step 4). So a
condition on \code{COUNT(*)} must be in HAVING.
\runsql{L2_having}
Both can appear together: first WHERE keeps only rejected applications, then GROUP BY counts them per candidate, then HAVING keeps candidates with at
least one.
\runsql{L2_where_vs_having}
\begin{intuition}
If a condition can be checked by looking at one row (status = 'rejected'), put it in WHERE: it shrinks the data before the expensive grouping. Only
conditions that need the whole group (``at least 3 applications'') belong in HAVING.
\end{intuition}

\section{Conditional aggregation: the pattern behind every rate}
\stuck{SQL conditional aggregation SUM CASE WHEN}{\yts{sql+conditional+aggregation+sum+case+when}}
To count only some rows inside each group, put a CASE inside the aggregate: \code{SUM(CASE WHEN status = 'offer' THEN 1 ELSE 0 END)} counts offers. The
average of a 0/1 column is the proportion of ones, i.e.\ a rate. This one idea gives offer rates, click-through rates, conversion rates, churn rates.
\runsql{L2_conditional_agg}
\runsql{L2_status_counts}
Job 207 had 3 applications and 1 offer: offer rate $1/3 = 0.33$. Note the \code{1.0} inside the CASE: with integer 1 the AVG would still work in SQLite,
but writing \code{1.0} protects you from integer division on every database. PostgreSQL also offers \code{COUNT(*) FILTER (WHERE status = 'offer')}.

\section{Joins: combining tables}
\stuck{SQL joins explained visually (inner, left, right, full)}{\yts{sql+joins+explained+visually+inner+left+right+full}}
\subsection{The idea, from zero}
The \code{jobs} table stores \code{company\_id = 2}, not the name ``PayWave''. To show the name, SQL must \emph{match} each job row with the company row
that has the same id. A join does exactly this: for every pair of rows (one from each table) it checks the \code{ON} condition and keeps the pairs for
which it is true.

\begin{example}[: a join by hand]
Left table L: (job 201, company 1), (job 203, company 2), (job 999, company 7). Right table R: (company 1, DataNest), (company 2, PayWave),
(company 6, GreenGrid). Condition: \code{L.company = R.company}.
\begin{itemize}
\item \textbf{INNER JOIN}: only matching pairs: (201, DataNest), (203, PayWave). Job 999 and GreenGrid disappear.
\item \textbf{LEFT JOIN}: all inner pairs, plus every left row with no match, padded with NULL: adds (999, NULL).
\item \textbf{RIGHT JOIN}: all inner pairs plus unmatched right rows: adds (NULL, GreenGrid).
\item \textbf{FULL OUTER JOIN}: inner pairs plus unmatched rows from both sides: 4 rows.
\item \textbf{CROSS JOIN}: every pair, no condition: $3 \times 3 = 9$ rows.
\end{itemize}
\end{example}

\begin{center}
\begin{tikzpicture}[scale=0.62, every node/.style={font=\scriptsize}]
\foreach \x/\lab/\fl/\fr/\fm in {0/INNER/white/white/accent!45, 4.6/LEFT/accent!45/white/accent!45, 9.2/RIGHT/white/accent!45/accent!45,
  13.8/FULL/accent!45/accent!45/accent!45}{
  \begin{scope}[shift={(\x,0)}]
  \fill[\fl] (0,0) circle (1); \fill[\fr] (1.2,0) circle (1);
  \begin{scope}\clip (0,0) circle (1); \fill[\fm] (1.2,0) circle (1);\end{scope}
  \draw (0,0) circle (1); \draw (1.2,0) circle (1);
  \node at (-0.45,0) {L}; \node at (1.65,0) {R};
  \node at (0.6,-1.4) {\textbf{\lab}};
  \end{scope}}
\end{tikzpicture}
\end{center}
The circles are a memory aid only. A join is \emph{not} a set intersection of rows: if one left row matches three right rows, it appears three times.
The real rule is ``every matching pair''.

\subsection{INNER JOIN}
\runsql{L2_inner_join}
\code{j} and \code{c} are \emph{table aliases}, short names so that \code{j.company\_id} says which table's column you mean (both tables have a
\code{company\_id} and a \code{city}, so you must qualify them). Joins chain: to describe each offer we need application $\to$ candidate, application
$\to$ job $\to$ company.
\runsql{L2_three_join}

\subsection{LEFT JOIN, and counting things that may be zero}
``How many jobs has each company posted?'' must include GreenGrid with 0. INNER JOIN would drop it, because it has no matching job. LEFT JOIN keeps every
company:
\runsql{L2_left_join}
Why \code{COUNT(j.job\_id)} and not \code{COUNT(*)}? For GreenGrid the LEFT JOIN produced one row whose job columns are all NULL. \code{COUNT(*)} counts that
row:
\runsql{L2_left_count_star}
GreenGrid shows 1 job it does not have. After a LEFT JOIN, always count a column from the right table.

\subsection{Anti-join: rows with no match}
``Candidates who never applied'': LEFT JOIN, then keep rows where the right side is NULL. (Chapter 47 shows the equivalent \code{NOT EXISTS}.)
\runsql{L2_anti_join}

\subsection{The WHERE-after-LEFT-JOIN trap}
We want every candidate with their offer job, or NULL if they have no offer. Putting the status filter in WHERE quietly turns the LEFT JOIN into an INNER
JOIN:
\runsql{L2_left_where_trap}
Only 4 candidates came back. The LEFT JOIN did produce rows for Kabir, Rohan, \dots, but with \code{a.status} NULL, and \code{WHERE a.status = 'offer'}
is UNKNOWN for them, so WHERE removed them. Put conditions on the right table \emph{inside ON}: they then decide which right rows match, not which result
rows survive.
\runsql{L2_left_on_filter}

\subsection{FULL OUTER JOIN}
Supply and demand by city: jobs per city vs candidates per city, keeping cities that appear on only one side.
\runsql{L2_full_join}
Delhi has 2 candidates and no jobs (unmet supply); Hyderabad has a job and no candidates (unmet demand): both are business insights that an inner join
would hide. \code{COALESCE(j.city, c.city)} picks whichever side has the city. MySQL has no FULL JOIN; emulate it with \code{LEFT JOIN ... UNION ...
RIGHT JOIN}.

\subsection{SELF JOIN: a table joined to itself}
\code{employees.manager\_id} points to another row of \code{employees}. Join the table to itself under two aliases, one playing ``employee'' and one
playing ``manager'':
\runsql{L2_self_join}
LEFT JOIN keeps Neha, who has no manager. The same pattern finds referral pairs among candidates:
\runsql{L2_referral_self}

\subsection{CROSS JOIN and non-equi joins}
A CROSS JOIN pairs every row with every row. It is useful to build a complete grid (every variant $\times$ every day), so that a later LEFT JOIN can show
zeros for missing combinations:
\runsql{L2_cross_join}
The ON condition does not have to be equality. A \emph{non-equi join} matches candidates to the jobs they are eligible for (enough experience, same city):
\runsql{L2_join_eligibility}
This is how a simple rule-based job recommender is written in SQL.

\subsection{How many rows does a join return? Fan-out}
\begin{itemize}
\item 1:1 key match: at most one row per left row.
\item 1:many (one job, many applications): each job row is repeated once per application.
\item many:many (joining on a non-unique column on both sides): the row count can explode ($m \times n$ per matching key).
\end{itemize}
Repetition is correct for row-level output, but \textbf{dangerous for sums}. Total advertised max-CTC across jobs:
\runsql{L2_fanout_truth}
The same SUM after joining applications:
\runsql{L2_fanout}
498 instead of 235: job 201 (max CTC 20) has 5 applications, so its 20 was added 5 times. Jobs with no application (210) vanished. This is \emph{fan-out}.
Fix: aggregate each table to the grain you need \emph{before} joining (Chapter 48 builds an ML feature table this way).

\begin{recap}[: joins]
Row count after a join = number of matching pairs (+ unmatched rows for outer joins). Before writing a join, ask ``what is one row of each table, and is
the join key unique on each side?'' Count rows before and after; if a SUM suddenly grows, suspect fan-out.
\end{recap}

\section{Set operations: UNION, UNION ALL, INTERSECT, EXCEPT}
\stuck{SQL UNION vs UNION ALL, INTERSECT, EXCEPT}{\yts{sql+union+vs+union+all+intersect+except}}
Joins put tables \emph{side by side}; set operations stack results \emph{on top of each other}. Both queries must return the same number of columns with
compatible types.
\begin{itemize}
\item \code{UNION}: stack and remove duplicates (needs a sort or hash, slower).
\item \code{UNION ALL}: stack and keep everything (faster; use it unless you need deduplication).
\item \code{INTERSECT}: rows in both. \code{EXCEPT} (Oracle: \code{MINUS}): rows in the first but not the second.
\end{itemize}
\runsql{L2_union}
\runsql{L2_union_all}
\runsql{L2_except}
Candidate cities with no jobs: Delhi, and NULL. Unlike \code{=}, set operations treat two NULLs as the same value, so the NULL city appears once.

% ======================================================================
\section{Interview questions: Level 2}
\stuck{SQL joins interview questions and row count puzzles}{\yts{sql+joins+interview+questions+row+count+duplicates+nulls}}

\begin{qa}{\classic Table A has values 1, 1, 2, NULL and table B has 1, 1, 3, NULL (one column \code{x}). How many rows do INNER, LEFT, RIGHT, FULL and
CROSS joins on \code{A.x = B.x} return?}
\oneline{INNER 4, LEFT 6, RIGHT 6, FULL 8, CROSS 16.}
\why A join returns every matching pair. The two 1s in A each match the two 1s in B: $2 \times 2 = 4$ pairs. 2 and 3 match nothing. NULL never equals
NULL, so the NULLs match nothing. LEFT adds A's unmatched rows (2 and NULL): $4 + 2 = 6$. RIGHT adds B's unmatched (3 and NULL): 6. FULL adds both: $4 + 2
+ 2 = 8$. CROSS: $4 \times 4 = 16$.
\worked Verified by running it:
\runsql{Q_join_counts}
\pushes \emph{What if A had three 1s?} INNER $= 3 \times 2 = 6$. \emph{How do you make NULLs match?} \code{ON A.x IS NOT DISTINCT FROM B.x} (or
\code{COALESCE(A.x, -1) = COALESCE(B.x, -1)} if $-1$ cannot occur).
\pitfall Answering ``INNER = 1 row because only the value 1 is common'' (thinking in sets instead of pairs), or matching the NULLs.
\end{qa}

\begin{qa}{\classic What is the difference between WHERE and HAVING? Can a query have both?}
\oneline{WHERE filters individual rows before grouping; HAVING filters groups after aggregation; yes, a query can and often does have both.}
\why Logical order: FROM $\to$ WHERE $\to$ GROUP BY $\to$ HAVING $\to$ SELECT. Aggregates do not exist at WHERE time, so \code{WHERE COUNT(*) > 2} is an
error. A row-level condition placed in HAVING works (on grouped columns) but wastes effort grouping rows that will be discarded.
\worked \code{L2\_where\_vs\_having}: WHERE keeps the 5 rejected applications, GROUP BY counts per candidate, HAVING keeps groups with at least one.
``Candidates with at least 3 applications'' needs \code{HAVING COUNT(*) >= 3}: 101 (4 rows, including the duplicate), 102, 103, 104, 105, 109.
\pushes \emph{Can HAVING be used without GROUP BY?} Yes; the whole table is then one group (\code{SELECT COUNT(*) FROM jobs HAVING COUNT(*) > 5}).
\emph{Where do you filter on a window function?} In an outer query, since windows are computed in SELECT, after HAVING.
\pitfall Saying ``HAVING is just WHERE for GROUP BY queries'' without the before/after distinction.
\end{qa}

\begin{qa}{\classic Find candidates who never applied to any job, in three different ways.}
\oneline{LEFT JOIN ... WHERE right.key IS NULL; NOT EXISTS (correlated subquery); NOT IN (subquery), the last only if the subquery cannot return NULL.}
\why All three express an anti-join. LEFT JOIN/IS NULL keeps unmatched left rows. NOT EXISTS checks, for each candidate, whether any application row
exists. NOT IN compares against a list; if that list contains NULL, \code{x NOT IN (..., NULL)} is never TRUE, and the query returns nothing.
\worked
\begin{lstlisting}[style=sql]
-- 1
SELECT c.cand_id, c.name FROM candidates c
LEFT JOIN applications a ON a.cand_id = c.cand_id WHERE a.app_id IS NULL;
-- 2
SELECT c.cand_id, c.name FROM candidates c
WHERE NOT EXISTS (SELECT 1 FROM applications a WHERE a.cand_id = c.cand_id);
-- 3 (safe here because applications.cand_id is NOT NULL)
SELECT cand_id, name FROM candidates
WHERE cand_id NOT IN (SELECT cand_id FROM applications);
\end{lstlisting}
All return Sana (106) and Priya (110). Third option with \code{EXCEPT}: \code{SELECT cand\_id FROM candidates EXCEPT SELECT cand\_id FROM applications}.
\pushes \emph{Which is fastest?} Modern optimisers turn NOT EXISTS and LEFT JOIN/IS NULL into the same anti-join plan; NOT EXISTS is the safest default.
Chapter 47 shows NOT IN returning 0 rows on \code{employees.manager\_id}, which contains a NULL.
\pitfall NOT IN against a nullable column.
\end{qa}

\begin{qa}{\classic UNION vs UNION ALL: which is faster and when must you use each?}
\oneline{UNION ALL simply concatenates and is faster; UNION also removes duplicate rows, which costs a sort or hash; use UNION only when duplicates are
wrong.}
\why Deduplication must compare all rows, $O(n\log n)$ sort or a hash table. When the two parts cannot overlap (e.g.\ January data UNION February data),
UNION ALL gives the same result for free. When stacking event logs from two sources where the same event might appear twice, UNION (or explicit
dedup) is needed.
\worked \code{L2\_union} returns 5 distinct cities; \code{L2\_union\_all} returns all 5 Bengaluru rows (2 companies + 3 candidates).
\pushes \emph{Do column names come from the first or second query?} The first. \emph{Can you ORDER BY inside each part?} Only the final result can be
ordered (wrap a part in a subquery if you need LIMIT inside it).
\pitfall Using UNION by habit on large tables, or mixing column orders between the two parts (types may line up silently but the meaning is swapped).
\end{qa}

\begin{qa}{Count jobs per company including companies with zero jobs. Why does \code{COUNT(*)} give the wrong answer?}
\oneline{\code{companies LEFT JOIN jobs ... GROUP BY company ... COUNT(j.job\_id)}; \code{COUNT(*)} counts the NULL-padded row a LEFT JOIN creates for
a company with no jobs, so it shows 1 instead of 0.}
\why A LEFT JOIN keeps an unmatched left row by padding right columns with NULL. That padded row is a real row of the result. \code{COUNT(j.job\_id)}
skips the NULL, giving 0.
\worked Outputs \code{L2\_left\_join} (GreenGrid 0) and \code{L2\_left\_count\_star} (GreenGrid 1).
\pushes \emph{And SUM of a right-table column for GreenGrid?} NULL, not 0; wrap with \code{COALESCE(SUM(...), 0)}.
\pitfall Using INNER JOIN, which drops GreenGrid entirely.
\end{qa}

\begin{qa}{Your LEFT JOIN returns fewer rows than the left table has. What happened?}
\oneline{A WHERE condition on a right-table column filtered out the NULL-padded rows, effectively converting the LEFT JOIN into an INNER JOIN; move the
condition into ON.}
\why For unmatched rows the right columns are NULL, so \code{WHERE right.col = 'x'} is UNKNOWN and the row is dropped. Conditions in ON restrict which
right rows can match, while still keeping every left row.
\worked \code{L2\_left\_where\_trap} returns 4 rows; \code{L2\_left\_on\_filter} returns all 10 candidates with offer job or NULL.
\pushes \emph{Is it ever correct to filter the right table in WHERE?} Yes, \code{WHERE right.key IS NULL} for an anti-join, or when you genuinely want
inner-join semantics (then write INNER JOIN, for clarity). \emph{What about conditions on the left table?} Those belong in WHERE; in ON of a LEFT JOIN
they would not remove left rows, which surprises people.
\pitfall Diagnosing ``missing rows'' as a data problem.
\end{qa}

\begin{qa}{After joining jobs to applications, the total of \code{max\_ctc\_lpa} doubled. Explain and fix.}
\oneline{Fan-out: each job row is repeated once per matching application, so its value is summed several times; aggregate to the job grain first, or sum
over the jobs table alone.}
\why Joining a one-side table (jobs) to a many-side table (applications) produces one row per application. Row-level questions (``list applications with
job info'') are fine; job-level sums are not.
\worked True sum 235 over 10 jobs vs 498 after the join over 23 rows (outputs \code{L2\_fanout\_truth}, \code{L2\_fanout}); job 201 alone contributes
$5 \times 20 = 100$ instead of 20.
\pushes \emph{How do you detect fan-out in a pipeline?} Assert row counts and key uniqueness before and after each join
(\code{COUNT(*) = COUNT(DISTINCT job\_id)}). \emph{Is \code{SUM(DISTINCT max\_ctc\_lpa)} a fix?} No: two different jobs with the same value would be
counted once.
\pitfall ``Fixing'' it with DISTINCT.
\end{qa}

\begin{qa}{Write a query for each company's offer rate (offers / applications), keeping companies with no applications.}
\oneline{LEFT JOIN companies $\to$ jobs $\to$ applications, then \code{1.0 * SUM(CASE WHEN status='offer' THEN 1 ELSE 0 END) / NULLIF(COUNT(app\_id), 0)}.}
\why Conditional aggregation counts offers; \code{1.0} avoids integer division; \code{NULLIF(x, 0)} turns a zero denominator into NULL, so the rate is
NULL (undefined) instead of a division-by-zero error. Both joins are LEFT so that GreenGrid survives. The duplicate application 22 is excluded inside
the ON.
\runsql{Q_company_offer_rate}
\pushes \emph{Is EduSpark (1 of 4) really better than PayWave (1 of 6)?} With these counts, no evidence; a Beta-smoothed rate or a confidence interval
(Chapter 38--39) is needed before ranking companies. \emph{Why put \code{app\_id <> 22} in ON, not WHERE?} In WHERE it would drop GreenGrid's NULL row.
\pitfall Dividing two integer counts.
\end{qa}

\begin{qa}{List each candidate's applications to jobs outside their own city. What happens to Sana?}
\oneline{Join applications to candidates and jobs and keep \code{ca.city <> j.city}; Sana would be dropped by the NULL comparison anyway, but she has no
applications.}
\runsql{Q_applied_outside_city}
\why Ten out-of-city applications: a relocation-willingness signal that a recommender can use. A candidate with NULL city would vanish from this list
because \code{NULL <> 'Mumbai'} is UNKNOWN; if such candidates existed you would have to decide whether ``unknown city'' counts as outside.
\pushes \emph{How would you turn this into a feature?} Per candidate: share of applications outside home city, via conditional aggregation.
\pitfall Not stating how NULL cities are treated.
\end{qa}

\begin{qa}{\deep The average of the department averages is 58.1, but the company-wide average salary is 53. Which is ``the average salary''?}
\oneline{53 is the average per employee; 58.1 is the average per department, which gives a one-person department (Leadership, 90) the same weight as a
four-person one; averages of averages are only equal to the overall average when groups have equal sizes.}
\runsql{Q_avg_of_avg}
\why The overall mean is $\sum_g n_g \bar x_g / \sum_g n_g$, a size-weighted mean of group means. The plain average of group means gives each group weight
1. Here department sizes (known salaries) are 4, 3, 1, 2 with means 45, 65, 90, 32.5: weighted $(180 + 195 + 90 + 65)/10 = 53$; unweighted $(45 + 65 +
90 + 32.5)/4 = 58.125$.
\pushes \emph{Where does this bite in analytics?} Averaging daily CTRs instead of total clicks / total impressions; averaging city-level conversion
rates for a national rate; computing ``average order value'' per store. Always ask what the unit of analysis is. \emph{Related paradox?} Simpson's
paradox: subgroup trends reversing in the aggregate (Chapter 48, A/B by day).
\pitfall Averaging pre-aggregated rates in dashboards.
\end{qa}

\begin{qa}{\deep How does a database physically execute a join? When is each method chosen?}
\oneline{Three main algorithms: nested-loop (good when one side is tiny or the inner side has an index), hash join (build a hash table on the smaller side,
probe with the larger; best for large equality joins), and sort-merge join (sort both on the key and walk them together; good when inputs are already
sorted or for range-like joins).}
\why Nested loop: for each outer row, find matches in the inner table, $O(n \cdot m)$ without an index, $O(n \log m)$ with one. Hash join: $O(n + m)$
time but needs memory for the hash table and only works for equality conditions. Sort-merge: $O(n\log n + m\log m)$, cheap if already sorted
(clustered index), handles inequality joins. The optimiser picks using table statistics (row counts, distinct values); \code{EXPLAIN} shows the choice.
\worked Joining 23 applications to 10 jobs on \code{job\_id}, the PK: a nested loop with an index lookup per application is trivial. Joining a billion
impression rows to 10 million jobs in Spark/Hive: a hash join, or a broadcast join if the jobs table fits in memory on every worker.
\pushes \emph{What is a broadcast join?} In distributed engines (Spark), the small table is copied to every machine so the big table need not be
shuffled across the network. \emph{Why can stale statistics slow a query?} The optimiser may think a table is tiny and choose a nested loop over millions
of rows.
\pitfall Assuming the written join order is the execution order.
\end{qa}
